%% Question bank chapter file - generated by bank_tools.py build (do not edit while the bank is still being built:
%% it is regenerated from the per-test data in Question Bank\data). Every question is one bankq block:
%%   \begin{bankq}[id=..., source={...}, paper=.., qno=.., type=..., subtopic={...}, answer={...}, answer_src=key|solved,
%%                 working={...}, difficulty=1-10, diff_parts={K S I C F}, diff_note={...},
%%                 group=<passage id or empty>, also_in={duplicate copies in other tests}]
%% difficulty = min(10, K+S+I+C+F-1): K concepts 1-3, S steps 1-4, I insight 0-2, C computation 0-2, F format 0-1
%%   ...question LaTeX, options, figures (\bfig / \bfigi -> img/)...
%%   \end{bankq}
%% Passages shared by several questions are bankpassage blocks; their questions carry group=<passage id>.
\documentclass[10pt]{article}
\usepackage[a4paper,margin=1.8cm]{geometry}
\usepackage[T1]{fontenc}
\usepackage{amsmath,amssymb,mathtools}
\usepackage{graphicx}
\usepackage[export]{adjustbox}
\usepackage{xcolor,needspace,keyval}
\usepackage{enumitem}
\usepackage{booktabs,array,tabularx}
\usepackage[hidelinks]{hyperref}
\setlength{\parindent}{0pt}
\setlength{\parskip}{3pt}
\makeatletter
\define@key{bq}{id}{\def\bq@id{#1}}\define@key{bq}{source}{\def\bq@src{#1}}\define@key{bq}{paper}{\def\bq@paper{#1}}
\define@key{bq}{qno}{\def\bq@qno{#1}}\define@key{bq}{type}{\def\bq@type{#1}}\define@key{bq}{subtopic}{\def\bq@sub{#1}}
\define@key{bq}{answer}{\def\bq@ans{#1}}\define@key{bq}{answer_src}{\def\bq@asrc{#1}}\define@key{bq}{working}{\def\bq@work{#1}}
\define@key{bq}{group}{\def\bq@group{#1}}\define@key{bq}{also_in}{\def\bq@also{#1}}
\define@key{bq}{difficulty}{\def\bq@diff{#1}}\define@key{bq}{diff_parts}{\def\bq@dparts{#1}}\define@key{bq}{diff_note}{\def\bq@dnote{#1}}
\newenvironment{bankq}[1][]{\def\bq@work{}\def\bq@also{}\def\bq@group{}\def\bq@diff{?}\def\bq@dparts{}\def\bq@dnote{}\setkeys{bq}{#1}%
  \par\needspace{4\baselineskip}\medskip\noindent{\small\textbf{\bq@id}\hfill\color{gray}[\bq@src, Paper \bq@paper, Q\bq@qno]
  \bq@type\ $\cdot$ \bq@sub\ $\cdot$ \textbf{D\bq@diff}}\par\nopagebreak\relax}
  {\par{\small\color{teal}Answer: \textbf{\bq@ans}\ifx\bq@work\@empty\else\ (solved: \bq@work)\fi
  \ifx\bq@also\@empty\else\par\color{gray}Also in: \bq@also\fi}\par}
\define@key{bp}{id}{\def\bp@id{#1}}\define@key{bp}{source}{\def\bp@src{#1}}\define@key{bp}{questions}{\def\bp@qs{#1}}
\newenvironment{bankpassage}[1][]{\setkeys{bp}{#1}\par\needspace{4\baselineskip}\bigskip\noindent
  {\small\textbf{Passage \bp@id}\hfill\color{gray}[\bp@src; questions \bp@qs]}\par\nopagebreak\relax\itshape}{\par}
\makeatother
\newcommand{\bfig}[2]{\par\smallskip{\centering\includegraphics[width=#1pt,max width=0.62\linewidth,max height=0.32\textheight]{img/#2}\par}\smallskip}
\newcommand{\bfigi}[2]{\includegraphics[width=#1pt,max width=\linewidth,max height=0.32\textheight,valign=c]{img/#2}}
\newcommand{\bmissing}[1]{\textit{[figure missing in the source conversion: \detokenize{#1}]}}
\newenvironment{qoptions}{\begin{itemize}[leftmargin=1.8em,itemsep=1pt,topsep=2pt,label={}]}{\end{itemize}}
\newcommand{\opt}[2]{\item[(#1)] #2}
\providecommand{\fwstart}{}\providecommand{\fwend}{}\providecommand{\qfullsplit}{}
\newcolumntype{L}[1]{>{\raggedright\arraybackslash}p{#1}}
\newcolumntype{C}[1]{>{\centering\arraybackslash}p{#1}}
\newcolumntype{R}[1]{>{\raggedleft\arraybackslash}p{#1}}
\providecommand{\sectionhdr}[1]{}\providecommand{\subsectionhdr}[1]{}


\begin{document}
{\Large\bfseries Aromatic --- Question Bank}\par\smallskip
{\small Chemistry chapter C23. Questions from all converted test papers, grouped by source series. Answers: printed key where it matches the paper, otherwise solved (answer\_src=solved, with working). Exact repeats across papers appear once, with also\_in listing the other copies.}\par\medskip\hrule

\section*{Score-I papers}
\begin{bankq}[id=SI12-P1-CHEM-S3-Q07, source={Score-I - Test 12 (09 Mar 2022)}, paper=1, qno=7, type=numerical, subtopic={Electrophilic substitution}, answer={6}, answer_src=key, difficulty=2, diff_parts={K1 S1 I0 C0 F1}, diff_note={concepts: activating and deactivating groups in electrophilic substitution | steps: classify the substituents; count rings more reactive than benzene}]
How many of the following compounds will show higher rate of electrophilic substitution reaction than benzene -
\bfig{220.1}{SI12_p1_chem_s3_q07_1.png}
\end{bankq}

\section*{Score Advanced papers}
\begin{bankq}[id=SA02-P1-CHEM-S3-Q05, source={Score Advanced - Test 02 (13 May 2021)}, paper=1, qno=5, type=numerical, subtopic={Directing effects}, answer={5}, answer_src=solved, working={Faster than p-bromonitrobenzene: fluoro compounds, extra nitro groups ortho or para to the halogen}, difficulty=7, diff_parts={K2 S2 I2 C1 F1}, diff_note={concepts: rate of aromatic nucleophilic substitution, directing effects | steps: activating groups; leaving group; count}]
How many of the following compounds given below have faster aromatic nucleophilic substitution\par
\bfig{308.0}{SA02_p1_chem_s3_q05_1.png}
\end{bankq}
\begin{bankq}[id=SA03-P2-CHEM-S2-Q08, source={Score Advanced - Test 03 (27 Jun 2021)}, paper=2, qno=8, type=single, subtopic={Directing effects}, answer={16}, answer_src=solved, working={m-nitrobenzaldehyde C7H5NO3}, difficulty=4, diff_parts={K2 S2 I1 C0 F0}, diff_note={concepts: nitration of benzaldehyde, directing effects | steps: Etard oxidation; meta nitration; atom count}]
\par
\bfig{119.9}{SA03_p2_chem_s2_q08_1.png}
Give the total number of atoms present in final product `B'.
\end{bankq}
\begin{bankpassage}[id=SA06-P1-CHEM-S2-G7-8, source={Score Advanced - Test 06 (04 Sep 2021)}, questions={7-8}]
Observe the following reaction sequence and give answers to questions given below :\par
\bfig{235.4}{SA06_p1_chem_s2_para7_8_1.png}
\end{bankpassage}
\begin{bankq}[id=SA06-P1-CHEM-S2-Q07, source={Score Advanced - Test 06 (04 Sep 2021)}, paper=1, qno=7, type=single, subtopic={Electrophilic substitution}, answer={A}, answer_src=key, difficulty=7, diff_parts={K3 S2 I2 C1 F0}, diff_note={concepts: Friedel-Crafts acylation and cyclisation, electrophilic substitution, benzene from hexane | steps: benzene from hexane; acylation with the anhydride; Clemmensen; cyclisation with AlCl3}, group=SA06-P1-CHEM-S2-G7-8]
Identify (Y) :
\begin{qoptions}
\opt{A}{\bfigi{68.9}{SA06_p1_chem_s2_q07_A.png}}
\opt{B}{\bfigi{47.3}{SA06_p1_chem_s2_q07_B.png}}
\opt{C}{\bfigi{72.3}{SA06_p1_chem_s2_q07_C.png}}
\opt{D}{\bfigi{59.2}{SA06_p1_chem_s2_q07_D.png}}
\end{qoptions}
\end{bankq}
\begin{bankq}[id=SA06-P1-CHEM-S2-Q08, source={Score Advanced - Test 06 (04 Sep 2021)}, paper=1, qno=8, type=single, subtopic={Electrophilic substitution}, answer={A}, answer_src=key, difficulty=5, diff_parts={K2 S1 I2 C1 F0}, diff_note={concepts: Hofmann exhaustive elimination from a ring amine, basicity | steps: methylation of the secondary amine; Hofmann elimination on each side}, group=SA06-P1-CHEM-S2-G7-8]
\bfig{169.5}{SA06_p1_chem_s2_q08_1.png}
\begin{qoptions}
\opt{A}{\bfigi{39.0}{SA06_p1_chem_s2_q08_A.png}}
\opt{B}{\bfigi{25.8}{SA06_p1_chem_s2_q08_B.png}}
\opt{C}{\bfigi{42.6}{SA06_p1_chem_s2_q08_C.png}}
\opt{D}{\bfigi{40.4}{SA06_p1_chem_s2_q08_D.png}}
\end{qoptions}
\end{bankq}
\begin{bankpassage}[id=SA07-P1-CHEM-S2-G7-8, source={Score Advanced - Test 07 (16 Apr 2023)}, questions={7-8}]
An organic compound A, C$_8$H$_4$O$_3$ , in dry benzene in the presence of anhydrous AlCl$_3$, gives compound B. The compound B on treatment with PCl$_5$ followed by reaction withH$_2$/Pd, BaSO$_4$ gives compound C, which on reaction with hydrazine gives a cyclised compound D(C$_{14}$H$_{10}$N$_2$).\par
\end{bankpassage}
\begin{bankq}[id=SA07-P1-CHEM-S2-Q07, source={Score Advanced - Test 07 (16 Apr 2023)}, paper=1, qno=7, type=single, subtopic={Electrophilic substitution}, answer={B}, answer_src=key, difficulty=5, diff_parts={K2 S1 I2 C1 F0}, diff_note={concepts: Friedel-Crafts with phthalic anhydride, electrophilic substitution | steps: A is phthalic anhydride; benzene gives 2-benzoylbenzoic acid}, group=SA07-P1-CHEM-S2-G7-8]
The compound `B' is
\begin{qoptions}
\opt{A}{\bfigi{73.5}{SA07_p1_chem_s2_q07_A.png}}
\opt{B}{\bfigi{67.7}{SA07_p1_chem_s2_q07_B.png}}
\opt{C}{\bfigi{95.6}{SA07_p1_chem_s2_q07_C.png}}
\opt{D}{\bfigi{93.1}{SA07_p1_chem_s2_q07_D.png}}
\end{qoptions}
\end{bankq}
\begin{bankq}[id=SA07-P1-CHEM-S2-Q08, source={Score Advanced - Test 07 (16 Apr 2023)}, paper=1, qno=8, type=single, subtopic={Electrophilic substitution}, answer={C}, answer_src=key, difficulty=8, diff_parts={K3 S2 I2 C2 F0}, diff_note={concepts: cyclised heterocycle from a keto aldehyde and hydrazine, electrophilic substitution, PCl5 gives the acid chloride | steps: PCl5 gives the acid chloride; Rosenmund gives the aldehyde; hydrazine cyclises}, group=SA07-P1-CHEM-S2-G7-8]
The compound `D' is
\begin{qoptions}
\opt{A}{\bfigi{40.2}{SA07_p1_chem_s2_q08_A.png}}
\opt{B}{\bfigi{54.4}{SA07_p1_chem_s2_q08_B.png}}
\opt{C}{\bfigi{43.1}{SA07_p1_chem_s2_q08_C.png}}
\opt{D}{\bfigi{43.0}{SA07_p1_chem_s2_q08_D.png}}
\end{qoptions}
\end{bankq}
\begin{bankq}[id=SA09-P1-CHEM-S3-Q06, source={Score Advanced - Test 09 (29 Apr 2023)}, paper=1, qno=6, type=single, subtopic={Electrophilic substitution}, answer={7}, answer_src=key, difficulty=4, diff_parts={K2 S1 I1 C1 F0}, diff_note={concepts: number of reactions with aromatic products, electrophilic substitution | steps: cyclisation and aromatisation routes}]
How many of the following will produce Aromatic product?\par
(A) $\text{CH}_3-\text{C}\equiv\text{CH} \xrightarrow[873\text{K}]{\text{Red hot iron Fe tube}}$ \quad (B) $\text{C}_6\text{H}_{14} \xrightarrow[\text{Cr}_2\text{O}_3/\Delta]{\text{Al}_2\text{O}_3}$ \quad (C) \bfigi{69.2}{SA09_p1_chem_s3_q06_1.png}\par
(D) Benzene $\xrightarrow[\text{Dark/Cold}]{\text{Cl}_2\text{,AlCl}_3\text{ excess}}$\par
\bfig{316.1}{SA09_p1_chem_s3_q06_2.png}
\end{bankq}
\begin{bankq}[id=SA12-P1-CHEM-S1-Q01, source={Score Advanced - Test 12 (11 May 2023)}, paper=1, qno=1, type=single, subtopic={Electrophilic substitution}, answer={D}, answer_src=key, difficulty=7, diff_parts={K3 S1 I2 C2 F0}, diff_note={concepts: reaction sequence leading to a compound used in quantitative analysis, electrophilic substitution, find the compound from the chlorination reaction | steps: find the compound from the chlorination reaction}, also_in={Score Advanced - Test 29 (01 May 2024) P2 Q1}]
Analyse the following reaction sequence.\par
\bfig{212.6}{SA12_p1_chem_s1_q01_1.png}
Which of the following compound(s) on reaction with Cl$_2$ can produce a compound X, That is used in quantitative estimation of ``Magnesium'' using carius method?
\begin{qoptions}
\opt{A}{X}
\opt{B}{W}
\opt{C}{Z}
\opt{D}{P}
\end{qoptions}
\end{bankq}
\begin{bankq}[id=SA13-P1-CHEM-S2-Q02, source={Score Advanced - Test 13 (13 May 2023)}, paper=1, qno=2, type=single, subtopic={Directing effects}, answer={5}, answer_src=key, difficulty=7, diff_parts={K3 S1 I2 C2 F0}, diff_note={concepts: compounds giving tribromo derivatives with bromine water, directing effects, activated rings with strongly donating groups | steps: activated rings with strongly donating groups}]
Consider the following compounds and count number of compounds which can produce tribromoderivatives on reaction with Br$_2$/H$_2$O.\par
\bfig{175.3}{SA13_p1_chem_s2_q02_1.png}
\bfig{221.8}{SA13_p1_chem_s2_q02_2.png}
\end{bankq}
\begin{bankq}[id=SA16-P1-CHEM-S4-Q09, source={Score Advanced - Test 16 (20 May 2023)}, paper=1, qno=9, type=numerical, subtopic={Directing effects}, answer={5}, answer_src=key, difficulty=5, diff_parts={K2 S1 I1 C1 F1}, diff_note={concepts: compounds nitrated more slowly than benzene, directing effects | steps: deactivating groups}]
In how many of the following compounds is the rate of nitration less than the rate of nitration of benzene.\par
\bfig{289.9}{SA16_p1_chem_s4_q09_1.png}
\end{bankq}
\begin{bankpassage}[id=SA16-P2-CHEM-S2-G7-8, source={Score Advanced - Test 16 (20 May 2023)}, questions={7-8}]
\par
\bfig{181.8}{SA16_p2_chem_s2_para7_8_1.png}
\end{bankpassage}
\begin{bankq}[id=SA16-P2-CHEM-S2-Q07, source={Score Advanced - Test 16 (20 May 2023)}, paper=2, qno=7, type=single, subtopic={Electrophilic substitution}, answer={B}, answer_src=key, difficulty=4, diff_parts={K2 S1 I1 C1 F0}, diff_note={concepts: incorrect statement about sulphonation of toluene, electrophilic substitution | steps: chlorosulphonation step}, group=SA16-P2-CHEM-S2-G7-8]
Choose the \textbf{INCORRECT} statement
\begin{qoptions}
\opt{A}{\textbf{[P]} \& \textbf{[Q]} are obtained from toluene by electrophilic substitution reaction}
\opt{B}{SO$_3$H will get substituted on benzene ring in first step and product formed is paramethylbenzenesulphonic acid and ortho-methylbenzenesulphonic acid.}
\opt{C}{\textbf{[P]} and \textbf{[Q]} are positional isomers}
\opt{D}{SO$_2$Cl will get substituted in first step}
\end{qoptions}
\end{bankq}
\begin{bankq}[id=SA16-P2-CHEM-S2-Q08, source={Score Advanced - Test 16 (20 May 2023)}, paper=2, qno=8, type=single, subtopic={Electrophilic substitution}, answer={A}, answer_src=key, difficulty=6, diff_parts={K3 S1 I2 C1 F0}, diff_note={concepts: saccharin synthesis, directing effects, which isomer is ortho | steps: which isomer is ortho; count oxygens and sp2 carbons}, group=SA16-P2-CHEM-S2-G7-8]
Choose the \textbf{CORRECT} Option
\begin{qoptions}
\opt{A}{\textbf{[P]} is the ortho isomer}
\opt{B}{\textbf{[Q]} is the ortho isomer}
\opt{C}{Number of oxygen atoms in \textbf{[S]} are 3}
\opt{D}{Number of sp$^2$ carbon atoms in saccharine are 6}
\end{qoptions}
\end{bankq}
\begin{bankq}[id=SA17-P2-CHEM-S1-Q04, source={Score Advanced - Test 17 (25 May 2023)}, paper=2, qno=4, type=single, subtopic={Electrophilic substitution}, answer={D}, answer_src=key, difficulty=7, diff_parts={K3 S1 I2 C2 F0}, diff_note={concepts: ordering of bromination rate constants, electrophilic substitution, electron donating groups favour the bromonium ion | steps: electron donating groups favour the bromonium ion}]
The rate constants for the bromination of several disubstituted stilbenes are given. Find the incorrect relation between rate constants for different X and Y?\par
The rate the table\par
\bfig{236.8}{SA17_p2_chem_s1_q04_1.png}
{\footnotesize\setlength{\tabcolsep}{4pt}\noindent\begin{tabular}{|c|c|c|}\hline
X & Y & K($\times10^3$ l/mol-min)\\\hline
OMe & Me & K$_1$\\\hline
OMe & Cl & K$_2$\\\hline
OMe & OMe & K$_3$\\\hline
OMe & NO$_2$ & K$_4$\\\hline
\end{tabular}}\par
\begin{qoptions}
\opt{A}{K$_3$ $>$ K$_1$ $>$ K$_2$}
\opt{B}{K$_1$ $>$ K$_2$ $>$ K$_4$}
\opt{C}{K$_3$ $>$ K$_2$ $>$ K$_4$}
\opt{D}{K$_1$ $>$ K$_3$ $>$ K$_4$}
\end{qoptions}
\end{bankq}
\begin{bankq}[id=SA19-P1-CHEM-S1-Q04, source={Score Advanced - Test 19 (04 Apr 2024)}, paper=1, qno=4, type=single, subtopic={Electrophilic substitution}, answer={B}, answer_src=key, difficulty=7, diff_parts={K3 S1 I2 C2 F0}, diff_note={concepts: sequence of Friedel Crafts bromination and nitration, electrophilic substitution, alkylation then cyclisation | steps: alkylation then cyclisation; directing effects}]
\par
\bfig{237.1}{SA19_p1_chem_s1_q04_1.png}
The product [M] is :
\begin{qoptions}
\opt{A}{\bfigi{46.7}{SA19_p1_chem_s1_q04_A.png}}
\opt{B}{\bfigi{65.9}{SA19_p1_chem_s1_q04_B.png}}
\opt{C}{\bfigi{62.2}{SA19_p1_chem_s1_q04_C.png}}
\opt{D}{\bfigi{65.1}{SA19_p1_chem_s1_q04_D.png}}
\end{qoptions}
\end{bankq}
\begin{bankq}[id=SA19-P1-CHEM-S3-Q08, source={Score Advanced - Test 19 (04 Apr 2024)}, paper=1, qno=8, type=numerical, subtopic={Directing effects}, answer={5}, answer_src=key, difficulty=10, diff_parts={K3 S3 I2 C2 F1}, diff_note={concepts: lone pairs in the product of a multistep aromatic synthesis, directing effects, chlorination | steps: chlorination; nitration; hydrolysis; reduction; acetylation}]
\par
\bfig{206.2}{SA19_p1_chem_s3_q08_1.png}
Find the number of lone pair of electron in product [T].
\end{bankq}
\begin{bankq}[id=SA22-P1-CHEM-S2-Q11, source={Score Advanced - Test 22 (10 Apr 2024)}, paper=1, qno=11, type=multi, subtopic={Directing effects}, answer={A,B,C}, answer_src=key, difficulty=8, diff_parts={K3 S1 I2 C2 F1}, diff_note={concepts: ortho and para nitration of toluene with given yields, directing effects, partial rate factors | steps: partial rate factors; relative rates}]
The nitration of toluene takes place at the ortho, para and meta position, giving ortho-, para-, and meta nitrotoluenes, respectively. Assume that the meta product is formed in negligible amount and that the ortho and para substitutions are irreversible first $-$ order reactions. Let the initial toluene concentration be [A]$_0$ and k$_1$ and k$_2$ be the ortho- and para $-$ rate constants, respectively. Choose the correct statements
\begin{qoptions}
\opt{A}{The rate of formation of the ortho substituted product is given by $k_1[A]\,e^{-(k_1+k_2)t}$, where t is the time.}
\opt{B}{If the rate constant for the ortho product formation is twice that for the para product formation, then [O-nitrotoluene]=2 [p-nitrotoluene], at any time}
\opt{C}{Nitration at the ortho $-$ and para $-$ positions is favorable compared to that at the meta position as the ortho and para $-$ intermediates are more stabilized than the meta $-$ intermediate}
\opt{D}{If the methyl group of toluene is replaced by tertiary butyl group, more ortho product is expected}
\end{qoptions}
\end{bankq}
\begin{bankpassage}[id=SA23-P1-CHEM-S3-G3-4, source={Score Advanced - Test 23 (14 Apr 2024)}, questions={3-4}]
Reaction scheme:\par
\bfig{207.7}{SA23_p1_chem_s3_para3_4_1.png}
\end{bankpassage}
\begin{bankq}[id=SA23-P1-CHEM-S3-Q03, source={Score Advanced - Test 23 (14 Apr 2024)}, paper=1, qno=3, type=numerical, subtopic={Directing effects}, answer={3}, answer_src=key, difficulty=8, diff_parts={K3 S1 I2 C2 F1}, diff_note={concepts: bromine atoms in the product of bromine water on a phenol derivative, directing effects, activating group | steps: activating group; ortho and para substitution}, group=SA23-P1-CHEM-S3-G3-4]
Total number of bromine atoms present in product when compound(Q) react with Br$_2$ $-$ H$_2$O ?
\end{bankq}
\begin{bankq}[id=SA23-P1-CHEM-S3-Q04, source={Score Advanced - Test 23 (14 Apr 2024)}, paper=1, qno=4, type=numerical, subtopic={Directing effects}, answer={4}, answer_src=key, difficulty=8, diff_parts={K3 S1 I2 C2 F1}, diff_note={concepts: nitrogen atoms in the final compound of a scheme, basicity, identify intermediates | steps: identify intermediates}, group=SA23-P1-CHEM-S3-G3-4]
Total number of Nitrogen atoms present in compound(U)?
\end{bankq}
\begin{bankq}[id=SA30-P1-CHEM-S2-Q09, source={Score Advanced - Test 30 (04 May 2024)}, paper=1, qno=9, type=multi, subtopic={Directing effects}, answer={B,D}, answer_src=key, difficulty=8, diff_parts={K3 S1 I2 C2 F1}, diff_note={concepts: standard problem in this chapter, directing effects, apply the relevant results | steps: apply the relevant results; compute}, also_in={Score Advanced - Test 31 (04 May 2024) P1 Q9}]
\bfig{131.3}{SA30_p1_chem_s2_q09_1.png}
P$_1$ \& P$_2$ will be distinguished by
\begin{qoptions}
\opt{A}{2, 4, D.N.P}
\opt{B}{alkaline Cu$^{2+}$ (tartarate)}
\opt{C}{NaHSO$_3$}
\opt{D}{Schiff's reagent}
\end{qoptions}
\end{bankq}
\begin{bankq}[id=SA30-P2-CHEM-S1-Q05, source={Score Advanced - Test 30 (04 May 2024)}, paper=2, qno=5, type=multi, subtopic={Electrophilic substitution}, answer={C,D}, answer_src=key, difficulty=8, diff_parts={K3 S1 I2 C2 F1}, diff_note={concepts: standard problem in this chapter, electrophilic substitution, apply the relevant results | steps: apply the relevant results; compute}, also_in={Score Advanced - Test 31 (04 May 2024) P2 Q5}]
Which of the following reaction/s is/are correctly represented with major products?
\begin{qoptions}
\opt{A}{\bfigi{132.8}{SA30_p2_chem_s1_q05_A.png}}
\opt{B}{\bfigi{170.4}{SA30_p2_chem_s1_q05_B.png}}
\opt{C}{\bfigi{126.8}{SA30_p2_chem_s1_q05_C.png}}
\opt{D}{\bfigi{183.8}{SA30_p2_chem_s1_q05_D.png}}
\end{qoptions}
\end{bankq}
\begin{bankq}[id=SA32-P2-CHEM-S2-Q08, source={Score Advanced - Test 32 (08 May 2024)}, paper=2, qno=8, type=multi, subtopic={Directing effects}, answer={A,B,D}, answer_src=key, difficulty=8, diff_parts={K3 S1 I2 C2 F1}, diff_note={concepts: standard problem in this chapter, directing effects, apply the relevant results | steps: apply the relevant results; compute}]
The reaction(s) leading to the formation of 1, 3, 5 trimethyl benzene is/are
\begin{qoptions}
\opt{A}{\bfigi{83.7}{SA32_p2_chem_s2_q08_A.png}}
\opt{B}{\bfigi{101.9}{SA32_p2_chem_s2_q08_B.png}}
\opt{C}{\bfigi{127.6}{SA32_p2_chem_s2_q08_C.png}}
\opt{D}{\bfigi{119.0}{SA32_p2_chem_s2_q08_D.png}}
\end{qoptions}
\end{bankq}
\begin{bankq}[id=SA33-P1-CHEM-S4-Q06, source={Score Advanced - Test 33 (12 May 2024)}, paper=1, qno=6, type=numerical, subtopic={Directing effects}, answer={39}, answer_src=key, difficulty=8, diff_parts={K3 S1 I2 C2 F1}, diff_note={concepts: standard problem in this chapter, directing effects, apply the relevant results | steps: apply the relevant results; compute}, also_in={Score Advanced - Test 37 (19 May 2024) P1 Q6}]
The correct order of the given reagents to convert benzene to m-Chlorobromobenzene is\par
(1) Sulfuric acid (conc.) and heat\par
(2) Cl$_2$ + FeCl$_3$ and heat\par
(3) NaNO$_3$ + H$_3$O$^+$, 0$^\circ$C\par
(4) H$_2$/Pt catalyst\par
(5) Mg in ether\par
(6) PBr$_3$\par
(7) H$_3$PO$_2$ (aqueous)\par
(8) HNO$_3$ (conc.) + H$_2$SO$_4$ (conc.) and heat\par
(9) Cu$_2$Br$_2$ +HBr\par
(10) (CH$_3$CO)$_2$ O + pyridine\par
Write the correct order in form of their number and express last two digits as answer.\par
[ For, e.g. if order is 1357, then your answer will be 57]
\end{bankq}
\begin{bankq}[id=SA34-P1-CHEM-S1-Q02, source={Score Advanced - Test 34 (12 May 2024)}, paper=1, qno=2, type=multi, subtopic={Electrophilic substitution}, answer={A,C,D}, answer_src=key, difficulty=8, diff_parts={K3 S1 I2 C2 F1}, diff_note={concepts: standard problem in this chapter, electrophilic substitution, apply the relevant results | steps: apply the relevant results; compute}, also_in={Score Advanced - Test 36 (19 May 2024) P1 Q2}]
Which of the following method(s) is/are not used to prepare p-bromo aniline as major product :
\begin{qoptions}
\opt{A}{\bfigi{72.8}{SA34_p1_chem_s1_q02_A.png}}
\opt{B}{\bfigi{172.0}{SA34_p1_chem_s1_q02_B.png}}
\opt{C}{\bfigi{119.6}{SA34_p1_chem_s1_q02_C.png}}
\opt{D}{\bfigi{189.2}{SA34_p1_chem_s1_q02_D.png}}
\end{qoptions}
\end{bankq}
\begin{bankq}[id=SA35-P2-CHEM-S1-Q04, source={Score Advanced - Test 35 (15 May 2024)}, paper=2, qno=4, type=single, subtopic={Electrophilic substitution}, answer={D}, answer_src=key, difficulty=7, diff_parts={K3 S1 I2 C2 F0}, diff_note={concepts: ordering of bromination rate constants, electrophilic substitution, electron donating groups favour the bromonium ion | steps: electron donating groups favour the bromonium ion}]
The rate\ \ constants for the bromination of several disubstituted stilbenes are given. Find the incorrect relation between rate constants for different X and Y?\par
The rate the table\par
\bfig{290.7}{SA35_p2_chem_s1_q04_1.png}

{\footnotesize\setlength{\tabcolsep}{4pt}\noindent\begin{tabular}{|c|c|c|}
\hline
X & Y & K($\times$10$^3$l/mol-min) \\ \hline
OMe & Me & K$_1$ \\ \hline
OMe & Cl & K$_2$ \\ \hline
OMe & OMe & K$_3$ \\ \hline
OMe & NO$_2$ & K$_4$ \\ \hline
\end{tabular}}\par
\begin{qoptions}
\opt{A}{K$_3$ $>$ K$_1$ $>$ K$_2$}
\opt{B}{K$_1$ $>$ K$_2$ $>$ K$_4$}
\opt{C}{K$_3$ $>$ K$_2$ $>$ K$_4$}
\opt{D}{K$_1$ $>$ K$_3$ $>$ K$_4$}
\end{qoptions}
\end{bankq}
\begin{bankq}[id=SA38-P1-CHEM-S4-Q05, source={Score Advanced - Test 38 (13 Apr 2025)}, paper=1, qno=5, type=numerical, subtopic={Electrophilic substitution}, answer={303}, answer_src=key, difficulty=8, diff_parts={K3 S1 I2 C2 F1}, diff_note={concepts: standard problem in this chapter, electrophilic substitution, apply the relevant results | steps: apply the relevant results; compute}]
In the following reaction sequence, the major product \textbf P is formed.\par
\bfig{126.7}{SA38_p1_chem_s4_q05_1.png}
1 mole of \textbf P react with $\dfrac13$ mole of glycerol in presence of an acid as catalyst gives product \textbf Q. Reaction of \textbf Q with excess NaOH followed by the treatment with CaCl$_2$ yields \textbf R, quantitatively.\par
Starting with $\dfrac13$ mole of \textbf Q, the amount of \textbf R produced in gm is \underline{\hspace{1.5cm}}\par
[Given, atomic weight: H = 1, C = 12, N = 14, O = 16, Na = 23, Cl = 35, Ca = 40]
\end{bankq}
\begin{bankq}[id=SA41-P1-CHEM-S1-Q03, source={Score Advanced - Test 41}, paper=1, qno=3, type=single, subtopic={Electrophilic substitution}, answer={C}, answer_src=key, difficulty=4, diff_parts={K2 S1 I1 C1 F0}, diff_note={concepts: standard problem in this chapter, electrophilic substitution | steps: apply the relevant results; compute}]
Sodium phenoxide when heated with CO$_2$ under pressure at 125$^\circ$C yields a product which on acetylation produces C.\par
\bfig{124.7}{SA41_p1_chem_s1_q03_R.png}
The major product C would be :
\begin{qoptions}
\opt{A}{\bfigi{64.1}{SA41_p1_chem_s1_q03_A2.png}}
\opt{B}{\bfigi{52.0}{SA41_p1_chem_s1_q03_B2.png}}
\opt{C}{\bfigi{62.9}{SA41_p1_chem_s1_q03_C2.png}}
\opt{D}{\bfigi{54.9}{SA41_p1_chem_s1_q03_D2.png}}
\end{qoptions}
\end{bankq}

\section*{Nurture SRG papers}
\begin{bankq}[id=NSRG03-P1-CHEM-S1-Q04, source={Nurture SRG - Test 03 (04 Feb 2024)}, paper=1, qno=4, type=single, subtopic={Directing effects}, answer={C}, answer_src=key, difficulty=5, diff_parts={K1 S2 I2 C1 F0}, diff_note={concepts: nitration of a fused heterocycle | steps: which ring is more activated; position of attack; major product}]
The major product of the following reaction is\par
\bfig{102.6}{NSRG03_p1_chem_s1_q04_1.png}
\begin{qoptions}
\opt{A}{\bfigi{67.1}{NSRG03_p1_chem_s1_q04_A.png}}
\opt{B}{\bfigi{102.7}{NSRG03_p1_chem_s1_q04_B.png}}
\opt{C}{\bfigi{91.8}{NSRG03_p1_chem_s1_q04_C.png}}
\opt{D}{\bfigi{103.5}{NSRG03_p1_chem_s1_q04_D.png}}
\end{qoptions}
\end{bankq}
\begin{bankq}[id=NSRG03-P1-CHEM-S1-Q06, source={Nurture SRG - Test 03 (04 Feb 2024)}, paper=1, qno=6, type=single, subtopic={Directing effects}, answer={C}, answer_src=key, difficulty=5, diff_parts={K1 S2 I2 C1 F0}, diff_note={concepts: product of a reaction on a fused ring | steps: sites of attack; the directing effect; major product}]
\bfig{159.9}{NSRG03_p1_chem_s1_q06_1.png}
\begin{qoptions}
\opt{A}{\bfigi{80.1}{NSRG03_p1_chem_s1_q06_A.png}}
\opt{B}{\bfigi{84.2}{NSRG03_p1_chem_s1_q06_B.png}}
\opt{C}{\bfigi{80.1}{NSRG03_p1_chem_s1_q06_C.png}}
\opt{D}{\bfigi{75.7}{NSRG03_p1_chem_s1_q06_D.png}}
\end{qoptions}
\end{bankq}
\begin{bankq}[id=NSRG03-P1-CHEM-S2-Q07, source={Nurture SRG - Test 03 (04 Feb 2024)}, paper=1, qno=7, type=multi, subtopic={Electrophilic substitution}, answer={A,B,C}, answer_src=solved, working={Strongly deactivated rings (dinitrobenzene, the ammonium substituted aldehyde, acetophenone) fail; pyrrole undergoes the reaction}, difficulty=4, diff_parts={K1 S2 I1 C0 F1}, diff_note={concepts: substrates failing in Friedel Crafts alkylation | steps: strongly deactivated rings; amines; heterocycles; select}]
Which of the following compound will not give Friedel-Crafts reaction ?
\begin{qoptions}
\opt{A}{\bfigi{96.2}{NSRG03_p1_chem_s2_q07_A.png}}
\opt{B}{\bfigi{58.3}{NSRG03_p1_chem_s2_q07_B.png}}
\opt{C}{\bfigi{75.1}{NSRG03_p1_chem_s2_q07_C.png}}
\opt{D}{\bfigi{32.6}{NSRG03_p1_chem_s2_q07_D.png}}
\end{qoptions}
\end{bankq}
\begin{bankq}[id=NSRG03-P1-CHEM-S3-Q01, source={Nurture SRG - Test 03 (04 Feb 2024)}, paper=1, qno=1, type=numerical, subtopic={Electrophilic substitution}, answer={5}, answer_src=solved, working={The tert-amyl cation is reached from 2-chloro-2-methylbutane, 3-chloro-2-methylbutane, 1-chloro-2-methylbutane, 1-chloro-3-methylbutane and neopentyl chloride (methyl shift)}, difficulty=5, diff_parts={K1 S2 I2 C0 F1}, diff_note={concepts: alkyl halides that give the same Friedel Crafts product | steps: target cation; rearrangements that give it; list the halides; count}]
Find the total number of possible alkyl chloride (R-Cl) in the following reaction.\par
\bfig{205.4}{NSRG03_p1_chem_s3_q01_1.png}
\end{bankq}
\begin{bankq}[id=NSRG03-P1-CHEM-S3-Q04, source={Nurture SRG - Test 03 (04 Feb 2024)}, paper=1, qno=4, type=numerical, subtopic={Directing effects}, answer={5}, answer_src=solved, working={Activating: OH, NH2, $\mathrm{CD_3}$, NHCOR and OCOR; the others deactivate}, difficulty=6, diff_parts={K1 S2 I2 C1 F1}, diff_note={concepts: substituents that activate the ring | steps: nature of each group; compare with benzene; count activating ones}]
Identify number of compounds which react faster than benzene in ArSE reaction ?\par
\bfig{208.0}{NSRG03_p1_chem_s3_q04_1.png}
\end{bankq}
\begin{bankq}[id=NSRG03-P2-CHEM-S3-Q06, source={Nurture SRG - Test 03 (04 Feb 2024)}, paper=2, qno=6, type=numerical, subtopic={Directing effects}, answer={6}, answer_src=solved, working={Toluene, anisole, phenyl acetate, pyrrole, naphthalene and the deuterated toluene are more reactive}, difficulty=5, diff_parts={K1 S2 I2 C0 F1}, diff_note={concepts: reactivity towards nitration | steps: effect of each substituent; compare with benzene; count}]
How many compounds are more reactive than benzene towards Nitration?\par
\bfig{245.2}{NSRG03_p2_chem_s3_q06_1.png}
\end{bankq}

\section*{Enthusiast Pre-Score papers}
\begin{bankq}[id=EPS01-P1-CHEM-S2-Q07, source={Enthusiast Pre-Score - Test 01 (04 Nov 2020)}, paper=1, qno=7, type=multi, subtopic={Directing effects}, answer={B,C,D}, answer_src=key, difficulty=3, diff_parts={K1 S1 I1 C0 F1}, diff_note={concepts: activated rings | steps: activating and deactivating groups; compare with benzene}]
Which of the following will react at rate faster than benzene during aromatic electrophilic substitution?
\begin{qoptions}
\opt{A}{\bfigi{29.3}{EPS01_p1_chem_s2_q07_A.png}}
\opt{B}{\bfigi{61.1}{EPS01_p1_chem_s2_q07_B.png}}
\opt{C}{\bfigi{40.2}{EPS01_p1_chem_s2_q07_C.png}}
\opt{D}{\bfigi{72.5}{EPS01_p1_chem_s2_q07_D.png}}
\end{qoptions}
\end{bankq}
\begin{bankq}[id=EPS01-P1-CHEM-S2-Q08, source={Enthusiast Pre-Score - Test 01 (04 Nov 2020)}, paper=1, qno=8, type=multi, subtopic={Electrophilic substitution}, answer={B,C}, answer_src=key, difficulty=5, diff_parts={K2 S2 I1 C0 F1}, diff_note={concepts: statements on electrophilic substitution, isotope effect | steps: isotope effect; effect of substituent; bromination of phenol; radical halogenation}]
Which are \textbf{CORRECT} statement(s) in the following?
\begin{qoptions}
\opt{A}{Rate of sulphonation of C$_6$H$_6$ and C$_6$D$_6$ are equal.}
\opt{B}{Rate of halogentation of\par\bfig{172.4}{EPS01_p1_chem_s2_q08_1.png}}
\opt{C}{\bfigi{29.0}{EPS01_p1_chem_s2_q08_2.png} on reaction with Br$_2$ / H$_2$O forms tribromoderivative.}
\opt{D}{\bfigi{31.6}{EPS01_p1_chem_s2_q08_3.png} on reaction with Cl$_2$ at 500$^\circ$C will give electrophilic aromatic substitution reaction.}
\end{qoptions}
\end{bankq}

\section*{TOAS papers}
\begin{bankq}[id=TOAS15-P2-CHEM-S1-Q02, source={TOAS - Test 15 (22 Aug 2022)}, paper=2, qno=2, type=multi, subtopic={Electrophilic substitution}, answer={A,B,D}, answer_src=solved, working={Acetone with conc. $\mathrm{H_2SO_4}$ and propyne over hot iron both trimerise to mesitylene; Clemmensen reduction of 1,3,5-triformylbenzene gives mesitylene; the triacetylbenzene haloform/soda-lime route gives benzene}, difficulty=4, diff_parts={K2 S2 I0 C0 F1}, diff_note={concepts: trimerisation, Clemmensen reduction | steps: trimerisation of acetone; trimerisation of propyne; Clemmensen reduction; the haloform route gives benzene}]
The reaction(s) leading to the formation of 1,3,5-trimethylbenzene is (are)
\begin{qoptions}
\opt{A}{\bfigi{84.5}{TOAS15_p2_chem_s1_q02_A.png}}
\opt{B}{Me--C$\equiv$C--H $\xrightarrow[\text{873 K}]{\text{heated iron tube}}$}
\opt{C}{\bfigi{122.5}{TOAS15_p2_chem_s1_q02_C.png}}
\opt{D}{\bfigi{137.0}{TOAS15_p2_chem_s1_q02_D.png}}
\end{qoptions}
\end{bankq}
\begin{bankq}[id=TOAS15-P2-CHEM-S2-Q04, source={TOAS - Test 15 (22 Aug 2022)}, paper=2, qno=4, type=numerical, subtopic={Directing effects}, answer={4.00}, answer_src=key, difficulty=2, diff_parts={K1 S1 I0 C0 F1}, diff_note={concepts: activating groups in electrophilic substitution | steps: classify the ten substituents as activating or deactivating; count}]
Number of compound show more rate of EAS than benezene.
\bfig{175.4}{TOAS15_p2_chem_s2_q04_1.png}
\end{bankq}
\begin{bankq}[id=TOAS16-P1-CHEM-S1-Q04, source={TOAS - Test 16 (04 Sep 2022)}, paper=1, qno=4, type=single, subtopic={Electrophilic substitution}, answer={B}, answer_src=key, difficulty=2, diff_parts={K2 S1 I0 C0 F0}, diff_note={concepts: Friedel-Crafts acylation, oxidation of a side chain | steps: Friedel-Crafts alkylation then permanganate gives the acid; acylation then Clemmensen does not}]
What reagents (conditions) are needed for the following reaction ?
\bfig{204.4}{TOAS16_p1_chem_s1_q04_1.png}
\begin{qoptions}
\opt{A}{reagent A : CH$_3$COCl/AlCl$_3$; reagent B : Sn, HCl, heat}
\opt{B}{reagent A : CH$_3$CH$_2$Cl/AlCl$_3$; reagent B : KMnO$_4$, heat}
\opt{C}{reagent A : CH$_3$COCl/AlCl$_3$; reagent B : H$_2$NNH$_2$, KOH, H$_2$O, heat}
\opt{D}{reagent A : HNO$_3$, H$_2$SO$_4$; reagent B : Sn, HCl, heat}
\end{qoptions}
\end{bankq}
\begin{bankq}[id=TOAS16-P1-CHEM-S2-Q10, source={TOAS - Test 16 (04 Sep 2022)}, paper=1, qno=10, type=multi, subtopic={Directing effects}, answer={A,C,D}, answer_src=key, difficulty=3, diff_parts={K1 S2 I0 C0 F1}, diff_note={concepts: substituent effects in electrophilic substitution | steps: nitro is meta directing and deactivating; chlorine is ortho para directing but deactivating; methoxy and methyl are activating}]
Which of the following statement is/are correct towards electrophilic substitution reaction when given groups are attached to benzene ring-
\begin{qoptions}
\opt{A}{--NO$_2$ is a metadirecting and deactivating group}
\opt{B}{--Cl is a metadirecting and deactivating group}
\opt{C}{--OMe is an ortho-paradirecting and activating group}
\opt{D}{--CH$_3$ is an ortho-paradirecting and activating group}
\end{qoptions}
\end{bankq}
\begin{bankpassage}[id=TOAS16-P1-CHEM-S3-G5-6, source={TOAS - Test 16 (04 Sep 2022)}, questions={5-6}]
Activating groups in benzene are more reactive and deactivating groups in benzene are less reactive.
\end{bankpassage}
\begin{bankq}[id=TOAS16-P1-CHEM-S3-Q05, source={TOAS - Test 16 (04 Sep 2022)}, paper=1, qno=5, type=numerical, subtopic={Directing effects}, answer={6}, answer_src=solved, working={Rings activated relative to benzene: toluene, phenyl benzoate (the O-bonded ring), biphenyl, phenyl acetate, dimethylaniline and phenol: 6}, difficulty=2, diff_parts={K1 S1 I0 C0 F1}, diff_note={concepts: activating groups in electrophilic substitution | steps: classify the nine compounds; count the activated rings}, group=TOAS16-P1-CHEM-S3-G5-6]
How many of the following compounds will show higher rate of electrophilic substitution reaction than benzene -
\bfig{213.0}{TOAS16_p1_chem_s3_q05_1.png}
\end{bankq}
\begin{bankq}[id=TOAS16-P1-CHEM-S3-Q06, source={TOAS - Test 16 (04 Sep 2022)}, paper=1, qno=6, type=numerical, subtopic={Directing effects}, answer={3.00}, answer_src=key, difficulty=4, diff_parts={K1 S2 I1 C0 F1}, diff_note={concepts: reactions of organic compounds | steps: predict the product of each of the ten reactions; compare with the given; count}, group=TOAS16-P1-CHEM-S3-G5-6]
How many reactions giving correct product?\par
\begin{center}\footnotesize\setlength{\tabcolsep}{2pt}
\begin{tabular}{@{}c@{\,}c@{}}
\bfigi{60.4}{TOAS16_p1_chem_s3_q06_1.png} & \bfigi{51.9}{TOAS16_p1_chem_s3_q06_2.png}\\
\bfigi{52.0}{TOAS16_p1_chem_s3_q06_4.png} & \bfigi{26.1}{TOAS16_p1_chem_s3_q06_3.png}\\
\bfigi{59.9}{TOAS16_p1_chem_s3_q06_6.png} & \bfigi{56.5}{TOAS16_p1_chem_s3_q06_5.png}\\
\bfigi{116.3}{TOAS16_p1_chem_s3_q06_8.png} & \bfigi{54.6}{TOAS16_p1_chem_s3_q06_7.png}\\
\bfigi{40.2}{TOAS16_p1_chem_s3_q06_10.png} & \bfigi{55.1}{TOAS16_p1_chem_s3_q06_9.png}\\
\end{tabular}\end{center}
\end{bankq}
\begin{bankq}[id=TOAS18-P1-CHEM-S2-Q10, source={TOAS - Test 18 (30 Oct 2022)}, paper=1, qno=10, type=single, subtopic={Electrophilic substitution}, answer={C}, answer_src=solved, working={P: cumene process gives phenol; Q: cyclopentenol dehydrates to cyclopentadiene and dichlorocarbene ring expansion gives chlorobenzene; R: Reimer-Tiemann acid reduced by zinc gives benzoic acid; S: acylation then permanganate gives terephthalic acid}, difficulty=3, diff_parts={K2 S1 I1 C0 F0}, diff_note={concepts: cumene process and carbene reactions and haloform, oxidation | steps: product of each sequence; match against the list}]
\begin{center}\footnotesize\setlength{\tabcolsep}{1pt}
\begin{tabular}{|L{0.6cm}L{3.8cm}|L{0.6cm}L{1.5cm}|}\hline
&List-I&&List-II\\\hline
(P)&\bfigi{122.4}{TOAS18_p1_chem_s2_q09_1.png}&(1)&\bfigi{33.8}{TOAS18_p1_chem_s2_q09_5.png}\\\hline
(Q)&\bfigi{124.2}{TOAS18_p1_chem_s2_q09_2.png}&(2)&\bfigi{27.9}{TOAS18_p1_chem_s2_q09_6.png}\\\hline
(R)&\bfigi{108.0}{TOAS18_p1_chem_s2_q09_3.png}&(3)&\bfigi{38.0}{TOAS18_p1_chem_s2_q09_7.png}\\\hline
(S)&\bfigi{144.5}{TOAS18_p1_chem_s2_q09_4.png}&(4)&\bfigi{26.5}{TOAS18_p1_chem_s2_q09_8.png}\\\hline
&&(5)&\bfigi{29.2}{TOAS18_p1_chem_s2_q09_9.png}\\\hline
\end{tabular}\end{center}
The correct option is :
\begin{qoptions}
\opt{A}{P $\to$ 1;Q $\to$ 2;R $\to$ 3;S $\to$ 4}
\opt{B}{P $\to$ 5;Q $\to$ 4;R $\to$ 3;S $\to$ 1}
\opt{C}{P $\to$ 2;Q $\to$ 5;R $\to$ 3;S $\to$ 1}
\opt{D}{P $\to$ 2;Q $\to$ 4;R $\to$ 5;S $\to$ 3}
\end{qoptions}
\end{bankq}
\begin{bankq}[id=TOAS18-P2-CHEM-S3-Q08, source={TOAS - Test 18 (30 Oct 2022)}, paper=2, qno=8, type=single, subtopic={Electrophilic substitution}, answer={4}, answer_src=key, difficulty=3, diff_parts={K1 S2 I1 C0 F0}, diff_note={concepts: bromination of aromatic compounds in schemes | steps: products of each scheme; final coupling product; count bromine atoms}]
Scheme 1 and 2 describe the conversion of P to Q and R to S, respectively. Scheme 3 describes the synthesis of T from Q and S. The total number of Br atoms in a molecule of T is \underline{\hspace{1.2cm}}
\par\noindent Scheme 1 :
\bfig{169.0}{TOAS18_p2_chem_s3_q08_1.png}
\par\noindent Scheme 2 :
\bfig{137.4}{TOAS18_p2_chem_s3_q08_2.png}
\par\noindent Scheme 3 :
\par\noindent\hfill$\mathbf S\xrightarrow[\text{(ii) }\mathbf Q]{\text{(i) NaOH}}\underset{\text{(major)}}{\mathbf T}$\hfill\null\par
\end{bankq}

\section*{TMAS papers}
\begin{bankq}[id=TMAS04-P1-CHEM-S3-Q05, source={TMAS - Test 04 (19 Nov 2023)}, paper=1, qno=5, type=numerical, subtopic={Directing effects}, answer={3.00}, answer_src=key, difficulty=4, diff_parts={K1 S1 I1 C1 F1}, diff_note={concepts: directing effects in nitration | steps: substituent effects; position}]
\bfig{147.7}{TMAS04_p1_chem_s3_q05_1.png}
At which position nitration mainly takes place ?
\end{bankq}

\section*{Star Review (Other) papers}
\begin{bankpassage}[id=STAR03-P1-CHEM-S3-G3-4, source={Star Review (Other) - Test 03 (04 Dec 2023)}, questions={3-4}]
Electrophilic aromatic substitution in phenols, the reaction that take place on the aromatic ring is electrophilic substitution reactions. The $-$OH group attached tot he benzene ring activates it towards electrophilic substitution. Also, it directs the incoming group to ortho and para positions in the ring as these positions become electron rich due to the mesomeric effect caused by $-$OH group.
\end{bankpassage}
\begin{bankq}[id=STAR03-P1-CHEM-S3-Q03, source={Star Review (Other) - Test 03 (04 Dec 2023)}, paper=1, qno=3, type=numerical, subtopic={Directing effects}, answer={5.00}, answer_src=key, difficulty=4, diff_parts={K1 S2 I1 C0 F1}, diff_note={concepts: Reimer Tiemann reaction and bromination of phenol | steps: position of the group on the ring; product of bromination; count}, group=STAR03-P1-CHEM-S3-G3-4]
\bfig{96.0}{STAR03_p1_chem_s3_q03_1.png}\par
Position of CHO group in A is $x$, no. of Br atoms in B is $y$ then $x+y=$ ?
\end{bankq}
\begin{bankq}[id=STAR03-P1-CHEM-S3-Q04, source={Star Review (Other) - Test 03 (04 Dec 2023)}, paper=1, qno=4, type=numerical, subtopic={Directing effects}, answer={6.00}, answer_src=key, difficulty=7, diff_parts={K1 S3 I2 C1 F1}, diff_note={concepts: isomers of a phenolic compound that rotate light | steps: degrees of unsaturation; phenol with a stereocentre; side chain options; ring positions; count}, group=STAR03-P1-CHEM-S3-G3-4]
An organic compound ($\mathrm{C_8H_{10}O_2}$) rotates plane-polarized light. It produces pink color with neutral $\mathrm{FeCl_3}$ solution. What is the total number of all the possible isomers for this compound?
\end{bankq}
\begin{bankpassage}[id=STAR03-P2-CHEM-S2-G7-8, source={Star Review (Other) - Test 03 (04 Dec 2023)}, questions={7-8}]
The $-\mathrm{NO_2}$ group in an aromatic ring deactivates the ortho and para positions for an electrophilic attack. When $-\mathrm{NO_2}$ group is present at ortho or para positions of a leaving group (Nucleofuge) it activates the ring for nucleophilic attack. The reduction of $-\mathrm{NO_2}$ group by metal in acid causes its reduction to $-\mathrm{NH_2}$ group and then the ring becomes strongly activated for a electrophilic attack. The strong activating group ($-\mathrm{NH_2}$) is converted into moderately activating group ($-\mathrm{NH{-}C(=O){-}CH_3}$, i.e.\ $-$NHAc) by it's acylation with $\mathrm{CH_3{-}CO{-}Cl}$. Deacylation is carried out by hydrolysis with $\mathrm{H_3O^+}$ or $\mathrm{OH^-}$. The ring alkylation by using $\mathrm{RX/AlX_3}$ is not possible in presence of $-\mathrm{NO_2}$ or $-\mathrm{NH_2}$ group but is possible in presence of $-$NHAc group.\par
\bfig{178.7}{STAR03_p2_chem_s2_para7_8_4.png}
\end{bankpassage}
\begin{bankq}[id=STAR03-P2-CHEM-S2-Q07, source={Star Review (Other) - Test 03 (04 Dec 2023)}, paper=2, qno=7, type=single, subtopic={Directing effects}, answer={B}, answer_src=key, difficulty=8, diff_parts={K3 S3 I2 C1 F0}, diff_note={concepts: multistep reaction with nitration, acylation, alkylation | steps: nitration; reduction; protection; alkylation; first product}, group=STAR03-P2-CHEM-S2-G7-8]
The product ($\mathrm{P_1}$) is :
\begin{qoptions}
\opt{A}{\bfigi{42.4}{STAR03_p2_chem_s2_q07_A.png}}
\opt{B}{\bfigi{59.3}{STAR03_p2_chem_s2_q07_B.png}}
\opt{C}{\bfigi{90.3}{STAR03_p2_chem_s2_q07_C.png}}
\opt{D}{\bfigi{63.0}{STAR03_p2_chem_s2_q07_D.png}}
\end{qoptions}
\end{bankq}
\begin{bankq}[id=STAR03-P2-CHEM-S2-Q08, source={Star Review (Other) - Test 03 (04 Dec 2023)}, paper=2, qno=8, type=single, subtopic={Directing effects}, answer={B}, answer_src=key, difficulty=8, diff_parts={K3 S3 I2 C1 F0}, diff_note={concepts: multistep reaction with nitration, acylation, alkylation | steps: nitration; reduction; protection; alkylation; second product}, group=STAR03-P2-CHEM-S2-G7-8]
The product ($\mathrm{P_2}$) is :
\begin{qoptions}
\opt{A}{\bfigi{37.7}{STAR03_p2_chem_s2_q08_A.png}}
\opt{B}{\bfigi{39.2}{STAR03_p2_chem_s2_q08_B.png}}
\opt{C}{\bfigi{41.4}{STAR03_p2_chem_s2_q08_C.png}}
\opt{D}{\bfigi{26.2}{STAR03_p2_chem_s2_q08_D.png}}
\end{qoptions}
\end{bankq}

\section*{Internal Test papers}
\begin{bankq}[id=INT01-P1-CHEM-S2-Q06, source={Internal Test - Test 01 (05 Aug 2020)}, paper=1, qno=6, type=multi, subtopic={Directing effects}, answer={A,B,C,D}, answer_src=key, difficulty=7, diff_parts={K1 S3 I2 C1 F1}, diff_note={concepts: reaction scheme of a phenolic acid | steps: identify A by decarboxylation; formaldehyde reactions; the resin; the other products; check each}]
Identify the correct statements in the below reaction scheme.\par
\bfig{269.9}{INT01_p1_chem_s2_pre_1.png}

\begin{qoptions}
\opt{A}{A $=$ \bfigi{60.9}{INT01_p1_chem_s2_pre_2.png}}
\opt{B}{D $=$ \bfigi{71.8}{INT01_p1_chem_s2_pre_B.png}}
\opt{C}{E $=$ \bfigi{77.0}{INT01_p1_chem_s2_pre_4.png}}
\opt{D}{F $=$ Bakelite}
\end{qoptions}
\end{bankq}

\section*{Nurture (General) papers}
\begin{bankq}[id=NUR01-P2-CHEM-S2-Q02, source={Nurture (General) - Test 01 (12 Jun 2021)}, paper=2, qno=2, type=multi, subtopic={Electrophilic substitution}, answer={A,D}, answer_src=solved, working={Excess benzene alkylation, and the NBS/aq KOH/PCC sequence; benzene with Cl2 gives BHC only with light}, difficulty=7, diff_parts={K2 S2 I2 C1 F1}, diff_note={concepts: four aromatic reactions, electrophilic substitution | steps: Friedel Crafts alkylation; chlorination conditions; NBS sequence}]
In given reaction, which one is/are correct according to major product.
\begin{qoptions}
\opt{A}{\bfigi{214.8}{NUR01_p2_chem_s2_q02_A.png}}
\opt{B}{\bfigi{158.3}{NUR01_p2_chem_s2_q02_2.png}}
\opt{C}{\bfigi{156.1}{NUR01_p2_chem_s2_q02_3.png}}
\opt{D}{\bfigi{100.3}{NUR01_p2_chem_s2_q02_4.png}}
\end{qoptions}
\end{bankq}
\begin{bankq}[id=NUR04-P1-CHEM-S2-Q07, source={Nurture (General) - Test 04 (11 Feb 2024)}, paper=1, qno=7, type=multi, subtopic={Electrophilic substitution}, answer={B,C,D}, answer_src=solved, working={Isobutyl chloride, isobutene and isobutanol give the tert-butyl cation; tert-butyl bromide with an alkoxide eliminates}, difficulty=7, diff_parts={K2 S2 I2 C1 F1}, diff_note={concepts: Friedel Crafts routes to tert-butylbenzene, electrophilic substitution | steps: elimination with a base; carbocation rearrangement; isobutene; alcohol with BF3}]
Among the following, reaction(s) which gives (give) tert-butyl benzene as the major product is(are)
\begin{qoptions}
\opt{A}{\bfigi{79.6}{NUR04_p1_chem_s2_q07_A.png}}
\opt{B}{\bfigi{78.7}{NUR04_p1_chem_s2_q07_B.png}}
\opt{C}{\bfigi{79.6}{NUR04_p1_chem_s2_q07_C.png}}
\opt{D}{\bfigi{78.7}{NUR04_p1_chem_s2_q07_D.png}}
\end{qoptions}
\end{bankq}
\end{document}
